\documentclass[11pt,a4paper]{article}
\usepackage[margin=27mm]{geometry}
\usepackage[T1]{fontenc}
\usepackage{lmodern}
\usepackage{amsmath,amssymb,amsthm,mathtools}
\usepackage{booktabs,longtable,array}
\usepackage{microtype}
\usepackage[hidelinks]{hyperref}
\usepackage{enumitem}
\setlist{itemsep=2pt,topsep=5pt}
\newtheorem{theorem}{Theorem}[section]
\newtheorem{lemma}[theorem]{Lemma}
\newtheorem{proposition}[theorem]{Proposition}
\newtheorem{corollary}[theorem]{Corollary}
\theoremstyle{remark}
\newtheorem{remark}[theorem]{Remark}
\newcommand{\fl}{\operatorname{fl}}
\newcommand{\epsin}{\varepsilon_{\mathrm{in}}}
\newcommand{\epsar}{\varepsilon_{\mathrm{ar}}}
\allowdisplaybreaks[2]
\setlength{\emergencystretch}{2em}
\title{A Numerical Certificate for the I.I.D. Lower Bound\\[3pt]
\large Correctness of the original FRLP evaluator}
\author{}
\date{}
\begin{document}
\maketitle
\begin{abstract}
This document supplies the implementation correctness arguments for the
fixed-coefficient factor-revealing LP in Section~3.3 and Appendix~A.1 of
\emph{Mean-only Prophet Inequalities}. It proves that the rate-tree
pruning rules, numerical table adjustments, interpolation, and explicit
arithmetic error allowances preserve all original LP inequalities.
The exact rational data and the frozen evaluator accompany the proof.
The result is a feasible certificate; no optimality of the LP solution
or of the chosen threshold is asserted.
\end{abstract}
\tableofcontents

\section{Certificate, scope, and notation}
The manuscript reference is \emph{Mean-only Prophet Inequalities}.
References such as Lemma~A.2 and Section~3.3 refer to that manuscript;
numbered results in this document have their own numbering.
The supplied rational threshold has $N=40$ pieces. Its initial value and
the certified objective are
\[
 \theta(0)=\frac{57507834299659}{77507834299659},\qquad
 \gamma=\frac{815733368877}{1099511627776}.
\]
The exact target comparison is
\[
10000\cdot815733368877-7419\cdot1099511627776
=56922299856>0.
\]
The proof below concerns feasibility of the original finite LP. The
paper's analytical transfer theorem gives the subsequent all-finite-$n$
consequence. The historical search for the threshold is not an input to
the correctness argument.

We retain the paper's symbols $b,\theta,c,p,u,\Delta,t_k,r_j$ and
$v_{k,i,j}$ with their original meanings. The integer arrays in the
program represent the exact rational values $v_{k,i,j}$ on a scale of
$2^{40}$. The state knots have denominator $2^{30}$.
The only interpolant used below, $\widetilde v_k$, is the ordinary
piecewise-affine interpolation of the stored layer-$k$ table.
It need not be concave. Concavity of the true function $V_t$, used to
relate the finite LP to the reward guarantee, is supplied by the paper.

The main files fixing this particular instance are:
\begin{center}
\small
\begin{tabular}{@{}p{0.52\linewidth}p{0.40\linewidth}@{}}
\toprule
File inside \texttt{certificate/} & Mathematical content\\
\midrule
\texttt{data/threshold\_slopes.json} & The exact rational drop rates $b_1,\ldots,b_{40}$\\
\texttt{data/grid1200.json} & The shared state knots $y_i=x_i$\\
\texttt{data/rates\_refined\_Q131093.json} & The finite rate endpoints; the last upper endpoint is $+\infty$\\
\texttt{data/coeff\_N640\_Q131093.h} & Binary64 inputs with proved rational enclosures\\
\texttt{verify.cpp} & The frozen backward evaluator\\
\texttt{validate\_inputs.py} & Exact-rational validation of the coefficient inputs\\
\bottomrule
\end{tabular}
\end{center}
The preflight script also checks that the text grid actually read by the
evaluator is identical to the validated JSON grid, and that the two
rational representations of the threshold agree.

% Mathematical core of the IID implementation supplement.
% Requires amsmath, amssymb, amsthm and environments lemma, theorem, proof.
% Arithmetic lemmas are supplied in a separate fragment.

\section{The original i.i.d. FRLP and its accelerated evaluation}
\label{sec:cert-iid-math}

This section proves that the accelerated backward evaluator constructs a
feasible solution of the fixed-coefficient FRLP in Section~3.3 of the paper.
The mathematical inputs are Lemmas~3.3, A.2, and A.3 of the paper.  The
additional arguments below justify the evaluation of that same finite LP:
they do not introduce stronger reward inequalities or replace any of its
rows.  Arithmetic lower enclosures are established separately below.  In
particular, every comparison used for pruning must use such an enclosure,
rather than a rounded estimate without a proved error bound.

\subsection{Fixed data and literal rows}
\label{sec:cert-iid-rows}

The certificate fixes $N=40$ positive nondecreasing drop rates, $N_t=640$
time intervals, $N_y=1200$ state subdivisions, and $N_q=131093$ rate
intervals.  Write $t_k=k/N_t$ and $\Delta=1/N_t$.  The state knots satisfy
\[
  0=y_0=x_0<y_1=x_1<\cdots<y_{N_y}=x_{N_y}=1,
\]
and the rate endpoints satisfy
\[
  0=r_0<r_1<\cdots<r_{N_q-1}<r_{N_q}=+\infty.
\]
All finite endpoints and drop rates are fixed rational inputs.  On each
state rectangle we use the diagonal from its lower-left to its upper-right
vertex, retaining only the part of the grid in $0\le x\le y\le1$.

Fix a time $t=t_k$, a state $0\le x<y\le1$, and the constant drop rate $b>0$
on $[t,t+\Delta]$.  As in Appendix~A.1, let
\[
 c=\max_{0\le s\le1}\{b(s)+\theta(s)\}-1>0,\qquad
 p=cx+y(1-\theta(t)-b),\qquad u=b(1-y+ty).
\]
For finite $r\ge0$, put
\[
 E(r)=e^{-\Delta r},\qquad
 I_r=\int_0^\Delta e^{-rs}\,ds
     =\begin{cases}(1-E(r))/r,&r>0,\\ \Delta,&r=0.\end{cases}
\]
The successor quantities in Lemma~A.2, viewed as functions of a finite
lower endpoint $r$, are
\begin{align}
 \varphi(r)&=1-y+y\bigl(t+I_r+E(r)(1-t-\Delta)\bigr),\label{eq:cert-iid-phi}\\
 y(r)&=\frac{yE(r)}{\varphi(r)},\label{eq:cert-iid-y}\\
 x(r)&=\min\!\left\{y(r),
   \frac{E(r)\bigl(cx+yb\Delta(r+e^{-r}-1)\bigr)}{c\varphi(r)}\right\}.
 \label{eq:cert-iid-x}
\end{align}
Here $\varphi(r)>0$ and $0\le x(r)\le y(r)\le1$.

For finite $r$, define the parameters of Lemma~A.3 by
\begin{align*}
 C(r)&=\min\{p\Delta,pI_r\},\\
 \alpha(r)&=ybE(r)\Delta^2+(-p)^+I_r\Delta,\\
 \beta(r)&=\Delta(u+ybI_r),\qquad z(r)=\beta(r)e^{-r}.
\end{align*}
Both $\alpha(r)$ and $\beta(r)$ are positive at a non-diagonal state.
The maximum of the two unrestricted bounds in Lemma~A.3 is
\begin{equation}
 \mathcal H(\alpha,\beta)=
 \begin{cases}
 \displaystyle \beta-\frac{\beta^2}{4\alpha},&\beta\le\alpha,\\[4pt]
 \displaystyle \alpha\left(\frac34+\frac12\log\frac\beta\alpha\right),
    &\beta>\alpha.
 \end{cases}
 \label{eq:cert-iid-free}
\end{equation}
Indeed, with $v=\beta/\alpha$, the difference between the quadratic and
logarithmic expressions, divided by $\alpha$, is
\[
 v-\frac{v^2}{4}-\frac34-\frac12\log v.
\]
It vanishes at $v=1$ and has derivative $-(v-1)^2/(2v)\le0$.
This proves the stated choice of branch.  When $z(r)>\alpha(r)/2$, write
\[
 J(r)=z(r)+\frac{\alpha(r)r}{2}
       -\frac{\alpha(r)^2}{24\bigl(z(r)-\alpha(r)/2\bigr)}.
\]
For this fixed time and state, suppress the indices $k,i,j$ in the cost coefficient. For the literal rate interval $[r_{h-1},r_h]$, let $a_h$ be the maximum of
\[
 \min\{p\Delta,0\},\qquad
 C(r_h)+\mathcal H\bigl(\alpha(r_h),\beta(r_h)\bigr),\qquad
 C(r_h)+J(r_h),
 \]
where the last term is included only when its guard holds, and the latter
two terms are omitted when $r_h=+\infty$.  Thus
$a_{N_q}=\min\{p\Delta,0\}$.  These are exactly the constants specified at
the end of Appendix~A.1, before numerical lower enclosures are taken.

Let $v_{k+1,i,j}$ be the already fixed next-layer table, and let
$\widetilde v_{k+1}(y,x)$ denote its ordinary continuous piecewise-affine
interpolant on the stated triangles.  The right-hand side of the literal
row indexed by $h$ is
\begin{equation}
 \mathcal R_h
 =a_h+\varphi(r_{h-1})\,
       \widetilde v_{k+1}\bigl(y(r_{h-1}),x(r_{h-1})\bigr).
 \label{eq:cert-iid-row}
\end{equation}
The current table entry must satisfy $v_{k,i,j}\le\mathcal R_h$ for every
$h$.  The code's zero-based leaf index $h-1$ uses \texttt{A(h-1)} for the
continuation and \texttt{B(h)} for the running cost
(\texttt{verify.cpp}, lines 63--96 and 134).

The evaluator may use a lower bound on $\mathcal R_h$ that is strictly
smaller than the literal expression.  In particular, omitting an
applicable lower bound from the maximum defining $a_h$, imposing a
stronger guard before using $J$, or subtracting a nonnegative allowance
cannot invalidate a candidate that is required to lie below the literal
row.  None of these operations changes the target LP.

\subsection{Lowering the numerical table}
\label{sec:cert-iid-lowering}

The monotonicity of the true functions $V_t$ does not by itself imply
monotonicity of an arbitrary numerical table.  We enforce the property
needed below by lowering the newly computed layer, in the following
order:
\begin{align*}
 &i=1,\ldots,N_y,\quad j=i-1,i-2,\ldots,0:
       &&v_{k,i,j}\leftarrow\min\{v_{k,i,j},v_{k,i,j+1}\};\\
 &i=1,\ldots,N_y,\quad j=0,\ldots,i-1:
       &&v_{k,i,j}\leftarrow\min\{v_{k,i,j},v_{k,i-1,j}\}.
\end{align*}
Diagonal entries remain zero.  These are the two passes in
\texttt{lower\_monotone} (lines 158--160).

\begin{lemma}[Feasibility and monotonicity after lowering]
\label{lem:cert-iid-lowering}
Suppose the next layer is fixed, the current layer is nonpositive with
zero diagonal, and every current literal row is satisfied.  The two
passes preserve those properties and yield
\[
 v_{k,i,j}\le v_{k,i,j+1},\qquad
 v_{k,i,j}\le v_{k,i-1,j}\qquad(0\le j<i\le N_y).
\]
Its interpolant is nonpositive, zero on $x=y$, nondecreasing in $x$, and
nonincreasing in $y$ throughout the triangular state domain.
\end{lemma}
\begin{proof}
Each row constraint for layer $k$ has a current entry only on its left-hand
side; its right-hand side uses the already fixed layer $k+1$.  Decreasing
a current entry therefore preserves every such constraint.  The passes
leave the diagonal unchanged and only decrease nonpositive entries.
They are applied before any earlier layer is constructed, so there are no
already constructed constraints from layer $k-1$ to preserve.

The descending order of $j$ in the first pass makes each row
nondecreasing in $j$.  In the second pass, row $i-1$ has already been
completed.  The new row $i$ is the pointwise minimum of its previous
nondecreasing row and the completed nondecreasing row $i-1$, with the
latter extended by a final zero.  Hence it remains nondecreasing.  Each
column inequality is imposed against the completed preceding row and
cannot be changed by a later update.  This proves both displayed
inequalities.

For the interpolation claim, consider a rectangle with coordinate widths
$d_y=y_{i+1}-y_i>0$ and $d_x=x_{j+1}-x_j>0$.  On its lower triangle with
vertices $(y_i,x_j)$, $(y_{i+1},x_j)$, $(y_{i+1},x_{j+1})$, the two
partial derivatives are
\[
 \partial_x\widetilde v_k
   =\frac{v_{k,i+1,j+1}-v_{k,i+1,j}}{d_x}\ge0,
 \qquad
 \partial_y\widetilde v_k
   =\frac{v_{k,i+1,j}-v_{k,i,j}}{d_y}\le0.
\]
On the upper triangle they are
\[
 \partial_x\widetilde v_k
   =\frac{v_{k,i,j+1}-v_{k,i,j}}{d_x}\ge0,
 \qquad
 \partial_y\widetilde v_k
   =\frac{v_{k,i+1,j+1}-v_{k,i,j+1}}{d_y}\le0.
\]
Only the lower triangle is present when $i=j$.  Continuity across edges
extends these monotonicities along any horizontal or vertical segment in
the domain.  Every interpolated value is a convex combination of
nonpositive entries, and interpolation along the diagonal uses only
zero diagonal entries.
\end{proof}

\subsection{A continuation bound for an entire rate subtree}
\label{sec:cert-iid-continuation}

A tree node consists of consecutive literal intervals.  Write $r_-$ for
its first lower endpoint and $r_+$ for its last upper endpoint; $r_-$ is
finite and $r_+$ may be infinite.  Every lower endpoint occurring in the
node lies in $[r_-,r_+]$.  Define
\[
 \rho(r)=\frac{x}{y}+\frac{b\Delta}{c}(r+e^{-r}-1),
 \qquad x(r)=y(r)\min\{1,\rho(r)\}.
\]

\begin{lemma}[Uniform continuation bound]
\label{lem:cert-iid-continuation}
Suppose $\widetilde v_{k+1}$ has the properties in
Lemma~\ref{lem:cert-iid-lowering}.  If $\rho(r_-)\ge1$, every continuation in
the node is zero.  Otherwise define
\[
 y_+=\begin{cases}y(r_+),&r_+<+\infty,\\0,&r_+=+\infty,\end{cases}
 \qquad x_-=y_+\rho(r_-).
\]
For every finite $r\in[r_-,r_+]$,
\begin{equation}
 \varphi(r)\widetilde v_{k+1}\bigl(y(r),x(r)\bigr)
 \ge
 \varphi(r_-)\widetilde v_{k+1}\bigl(y(r_-),x_-\bigr).
 \label{eq:cert-iid-node-continuation}
\end{equation}
The bound remains valid if $y_+$ is replaced by any certified value in
$[0,y_+]$, in particular by zero.
\end{lemma}
\begin{proof}
The functions $E(r)$ and $I_r$ are nonincreasing.  Since
$1-t-\Delta\ge0$, equation~\eqref{eq:cert-iid-phi} shows that
$\varphi$ is nonincreasing.  Moreover,
\[
 \frac{\varphi(r)}{E(r)}
 =(1-y)e^{\Delta r}
  +y\left(te^{\Delta r}+\int_0^\Delta e^{rs}\,ds+1-t-\Delta\right)
\]
is nondecreasing, so $y(r)=y/(\varphi(r)/E(r))$ is nonincreasing.
Finally $\rho'(r)=b\Delta(1-e^{-r})/c\ge0$.
If $\rho(r_-)\ge1$, it follows that $x(r)=y(r)$, where the interpolant is
zero.

Otherwise $0\le \rho(r_-)<1$.  For every finite $r$ under consideration,
$y(r)\ge y_+$ and $\min\{1,\rho(r)\}\ge \rho(r_-)$; hence
$x(r)\ge x_-$.  Also $y(r)\le y(r_-)$, and
$0\le x_-\le y_+\le y(r_-)$, so the comparison state belongs to the
domain.  Coordinate monotonicity gives
\[
 \widetilde v_{k+1}\bigl(y(r),x(r)\bigr)
 \ge \widetilde v_{k+1}\bigl(y(r_-),x_-\bigr).
\]
Multiplying by $\varphi(r)>0$, and then using
$\varphi(r)\le\varphi(r_-)$ and the nonpositivity of the right-hand
interpolated value, proves~\eqref{eq:cert-iid-node-continuation}.
Replacing $y_+$ by a smaller nonnegative value only decreases $x_-$ and
therefore weakens this lower bound.
\end{proof}

This is the exact-arithmetic argument behind \texttt{rectA}
(lines 98--106).  The implementation uses zero for the far endpoint
unless its arithmetic guard permits an accurately enclosed division.
Its coordinate perturbations and multiplier rounding are justified by
the arithmetic enclosures below; they are not assumptions of exact
monotonicity for rounded expressions.

\subsection{A running-cost bound for an entire rate subtree}
\label{sec:cert-iid-cost}

\begin{lemma}[Monotonicity of the free cost bound]
\label{lem:cert-iid-free}
For fixed $t,x,y,b$, the functions $C(r)$, $\alpha(r)$, and $\beta(r)$
are nonincreasing in finite $r\ge0$.  The function
$\mathcal H(\alpha,\beta)$ is nondecreasing in each positive argument.
Consequently, if $r_+<+\infty$, then every literal cost in the node is
at least
\begin{equation}
 C(r_+)+\mathcal H\bigl(\alpha(r_+),\beta(r_+)\bigr).
 \label{eq:cert-iid-node-free}
\end{equation}
Every literal cost, including the unbounded leaf, is also at least
$\min\{p\Delta,0\}$.
\end{lemma}
\begin{proof}
If $p<0$, then $C(r)=p\Delta$ because $I_r\le\Delta$; if $p\ge0$,
then $C(r)=pI_r$.  Thus $C$ is nonincreasing.  The formulas for $\alpha$
and $\beta$ are nonnegative linear combinations of $E(r)$, $I_r$, and
constants, so the same conclusion holds for them.  For $\beta\le\alpha$,
\[
 \partial_\alpha \mathcal H=\frac{\beta^2}{4\alpha^2}>0,
 \qquad \partial_\beta \mathcal H=1-\frac{\beta}{2\alpha}\ge\frac12.
\]
For $\beta\ge\alpha$,
\[
 \partial_\alpha \mathcal H=\frac14+\frac12\log\frac\beta\alpha>0,
 \qquad \partial_\beta \mathcal H=\frac\alpha{2\beta}>0.
\]
The formulas and their first derivatives agree at $\beta=\alpha$, which
proves coordinate monotonicity on the whole positive quadrant.  Each
finite leaf in the node has upper endpoint at most $r_+$, and its
literal cost includes this free bound evaluated at its own upper
endpoint.  The claimed inequalities follow.
\end{proof}

\begin{lemma}[Uniform shifted cost bound]
\label{lem:cert-iid-shifted}
Suppose $r_+<+\infty$ and $z(r_+)>\alpha(r_-)$.  Then every literal
cost in the node is at least
\begin{equation}
 C(r_+)+z(r_+)+\frac{\alpha(r_+)r_-}{2}
 -\frac{\alpha(r_-)^2}
 {24\bigl(z(r_+)-\alpha(r_-)/2\bigr)}.
 \label{eq:cert-iid-node-shifted}
\end{equation}
\end{lemma}
\begin{proof}
Both factors in $z(r)=\beta(r)e^{-r}$ are nonnegative and
nonincreasing, so $z$ is nonincreasing.  At every finite upper endpoint
$r\in[r_-,r_+]$ of a descendant leaf,
\[
 z(r)\ge z(r_+)>\alpha(r_-)\ge\alpha(r)>\alpha(r)/2.
\]
Thus its literal cost includes the shifted bound $C(r)+J(r)$.
The positive terms satisfy
\[
 C(r)\ge C(r_+),\qquad z(r)\ge z(r_+),\qquad
 \alpha(r)r\ge\alpha(r_+)r_-.
\]
For the subtracted term, its denominator is positive and
\[
 \frac{\alpha(r)^2}{z(r)-\alpha(r)/2}
 \le
 \frac{\alpha(r_-)^2}{z(r_+)-\alpha(r_-)/2}.
\]
Combining these inequalities proves the result.  No monotonicity of the
shifted expression $J(r)$ itself is asserted or required.
\end{proof}

The maximum of the available bounds in the preceding two lemmas is
therefore a valid running-cost lower bound for the node.  For a node
with $r_+=+\infty$ we retain only $\min\{p\Delta,0\}$.  This is the
argument used by \texttt{Bnode} (lines 108--124).  Its stronger numerical
guards, including $r_+\le64$, only suppress optional bounds.  The
implemented free bound uses a conservative logarithm evaluation in
\texttt{log\_lower}; no polynomial replacement of $\mathcal H$ is used or
needed for this certificate.

\subsection{Complete coverage and backward feasibility}
\label{sec:cert-iid-backward}

For each node, add a certified lower enclosure of
\eqref{eq:cert-iid-node-continuation} to a certified lower enclosure of
an applicable running-cost bound.  Call the resulting quantity
$\underline{\mathcal R}_{\rm node}$.  By the preceding lemmas,
\begin{equation}
 \underline{\mathcal R}_{\rm node}\le\mathcal R_h
 \quad\text{for every literal leaf }h\text{ in that node}.
 \label{eq:cert-iid-node-enclosure}
\end{equation}
For an individually evaluated leaf, write
$\underline{\mathcal R}_h\le\mathcal R_h$ for the certified enclosure
returned by its continuation and cost evaluations.

The default run subtracts $2^{-28}$ from each candidate.
It also subtracts a nonnegative quantity of the form
\[
  2^{-40}\left\lceil 2^{24}d_h\right\rceil,
  \qquad 0\le d_h\le y-x\le1,
\]
where the last inequalities are enforced by clamping in
\texttt{allowance} (lines 126--132).  Its detailed formula is only a
candidate-selection policy.  For feasibility, the only fact needed is
that the subtracted quantity is nonnegative.  In particular,
\[
 \underline{\mathcal R}_h-2^{-28}-2^{-40}\lceil 2^{24}d_h\rceil\le\mathcal R_h.
\]
No assertion about a telescoping potential or approximation to an exact
potential decrement is used.

\begin{theorem}[Feasibility of the accelerated backward evaluator]
\label{thm:cert-iid-backward}
Assume the arithmetic routines return the lower enclosures specified
above.  Starting from the zero terminal table, the evaluator's recursive
search and subsequent lowering passes produce a feasible solution of
the original FRLP.
\end{theorem}
\begin{proof}
Fix a state and a previously fixed next layer.  The algorithm maintains a
nonpositive incumbent $w$, initialized by the minimum of zero and
conservative candidate values.  Every update takes a minimum with
$\underline{\mathcal R}_h-2^{-28}-2^{-40}\lceil 2^{24}d_h\rceil$, so $w$ can only decrease.
The initial hint, the unbounded leaf, and the ternary-search proposals
supply candidates only; they do not remove any leaf from the subsequent
search.

A node is pruned only if
$\underline{\mathcal R}_{\rm node}\ge w$.  At that moment,
\eqref{eq:cert-iid-node-enclosure} implies $w\le\mathcal R_h$ for every
leaf in that node.  Later decreases of $w$ preserve all these inequalities.
If a node is not pruned and is a singleton, its leaf is evaluated and the
incumbent is updated; the resulting value is at most that leaf's literal
row.  Otherwise the node is partitioned into two disjoint nonempty
consecutive sets of leaves and both children are searched.

The root contains all $N_q$ leaves, including the last unbounded
interval.  Recursive splitting strictly decreases the number of leaves,
so the search terminates.  Its terminal nodes form a disjoint cover of
the root.  Every terminal node is either pruned with its inequalities
proved, or explicitly evaluated.  Consequently the final incumbent
satisfies $w\le\mathcal R_h$ for every $h$.  This argument also covers an
immediate pruning of the root.  The source's coverage counter, which
adds the number of leaves at a pruned node and one at an evaluated
singleton, checks this complete cover; the counter alone is not a
substitute for~\eqref{eq:cert-iid-node-enclosure}.

Apply this search independently to each non-diagonal state.  Its outputs
are nonpositive and satisfy every literal row against the fixed next
layer; the diagonal is left at zero.  Lemma~\ref{lem:cert-iid-lowering}
then preserves feasibility and supplies the interpolant properties
needed at the next backward step.  The zero terminal table has these
properties.  Backward induction from $k=N_t-1$ to $k=0$ proves the claim.
\end{proof}

The search just analyzed is \texttt{solve}/\texttt{search}
(lines 135--155); the layer update followed by lowering is in
lines 190--201.  In particular, no assumption of unimodality is needed
for the ternary-search proposals, and no value obtained from an ancestor
node is substituted for a literal leaf row.

Finally, the initial table entry is an integer multiple of $2^{-40}$,
and the reported objective is
\[
 \gamma
  =2^{-40}\lfloor2^{40}\theta(0)\rfloor+v_{0,N_y,0}
  \le\theta(0)+v_{0,N_y,0}.
\]
The FRLP implication proved in Section~3.3 of the paper, using the
concavity of the true $V_t$ and Lemma~3.3, therefore gives
\[
 L(b,q)\ge\theta(0)+V_0(1,0)
       \ge\theta(0)+v_{0,N_y,0}
       \ge\gamma
\]
for every admissible $q$.  The numerical interpolant itself need not be
concave.  The all-finite-$n$ consequence is supplied by the paper's
separate transfer theorem; it is not an additional assertion of this
backward evaluator.


\section{Certified arithmetic for the frozen evaluator}
\label{sec:arithmetic}
The numerical analysis applies to the exact source file identified in
\texttt{PAPER\_REFERENCE.json}. The numerical error constants are
\[
 \varepsilon_{\mathrm{in}}=2^{-48},\qquad
 \varepsilon_{\mathrm{ar}}=2^{-32}
   =65536\varepsilon_{\mathrm{in}},\qquad
 \mathfrak u=2^{-53}.
\]
Here $\mathfrak u$ is the unit roundoff of binary64; the paper's symbol
$u=b(1-y+ty)$ retains its mathematical meaning. The indexed LP
coefficients $\eta_{k,i,j,h,i',j'}$ also retain their original meaning;
they are unrelated to these numerical error constants.

% Numerical continuation proof for the frozen IID verifier.
% Intended to follow the exact-arithmetic pruning and monotonicity lemmas.
\subsection{Certified evaluation of the continuation terms}
\label{sec:certificate-continuation-roundoff}

This subsection concerns the frozen instance with $N_t=640$, state knots
that are integer multiples of $2^{-30}$, and finite rate endpoints at most
$8192$. It justifies the continuation calculations in
\texttt{verify.cpp}, including the tests that return a zero continuation.
Recall the rounding allowances
\[
 \varepsilon_{\mathrm{in}}=2^{-48},\qquad
 \varepsilon_{\mathrm{ar}}=2^{-32}
       =65536\varepsilon_{\mathrm{in}},\qquad
 \mathfrak u=2^{-53}=\varepsilon_{\mathrm{in}}/32.
\]
The code calls the first two quantities \texttt{eta} and \texttt{eps}.
All floating-point operations below are binary64 operations in
round-to-nearest mode with gradual underflow (no flush-to-zero or
denormals-are-zero mode), without contraction or reassociation. We write
$\operatorname{fl}$ for the result of one such operation. The usual
bound, also valid when its result is subnormal, is
\begin{equation}
 |\operatorname{fl}(z)-z|
 \le \mathfrak u|z|+2^{-1075}.
 \label{eq:certificate-fp-model}
\end{equation}
None of the operations considered below overflows. The displayed bounds
retain ample strict slack for the additive term in
\eqref{eq:certificate-fp-model}; in particular, a positive state coordinate
satisfies $y\ge 2^{-30}$, and
$2^{-1075}<2^{-990}\varepsilon_{\mathrm{in}}y$.

\paragraph{Exact quantities and validated inputs.}
Put $\Delta=1/640$ and $t=k\Delta$, where $0\le k<640$.
At a nonterminal state, $0\le x<y\le1$.
The divisions that form $x$ and $y$ from the integer knots and $2^{30}$
are exact. Write $b$ for the exact drop rate on this time interval and
$D=1-\theta(t)$. The input validator establishes
\[
 \tfrac14<b<2,\qquad \tfrac{17}{20}<c<\tfrac9{10},\qquad
 0\le D\le1.
\]
For a finite exact endpoint $r\in[0,8192]$, define
\[
 E(r)=e^{-\Delta r},\quad
 I_r=\int_0^\Delta e^{-rs}\,ds.
\]
Thus $0<E(r)\le1$, $0<I_r\le\Delta$, and $0\le (r+e^{-r}-1)\le8192$.
The frozen values of $b,c,D,r,E(r),I_r,(r+e^{-r}-1)$ differ from these exact
quantities by less than $\varepsilon_{\mathrm{in}}$ in absolute value.
The source entry \texttt{vv[h]} is the frozen approximation of
$r+e^{-r}-1$. Both arrays \texttt{bs} and \texttt{bl} approximate the same exact $b$.
The frozen $\widehat\Delta$ satisfies the stronger relative bound
\[
 |\widehat\Delta-\Delta|<\Delta\varepsilon_{\mathrm{in}}.
\]
The computed $\widehat t=\operatorname{fl}(k\widehat\Delta)$ consequently
satisfies
\begin{equation}
 |\widehat t-t|
 \le k\Delta\varepsilon_{\mathrm{in}}
   +\mathfrak u k\Delta(1+\varepsilon_{\mathrm{in}})+2^{-1075}
 <2\varepsilon_{\mathrm{in}}.
 \label{eq:certificate-time-error}
\end{equation}
Input validation uses rational enclosures; these input-error assumptions
therefore do not assume accuracy of a platform exponential routine.

Recall the exact continuation quantities
\begin{align*}
 \varphi(r)&=1-y+y\bigl[t+I_r+E(r)(1-t-\Delta)\bigr],\\
 y(r)&=\frac{yE(r)}{\varphi(r)},\qquad
 \rho(r)=\frac{x}{y}+\frac{b\Delta (r+e^{-r}-1)}c,\\
 \widetilde x(r)&=y(r)\rho(r),\qquad
 x(r)=\min\{y(r),\widetilde x(r)\}.
\end{align*}
We will use
\begin{equation}
 E(r)\le\varphi(r)\le1,
 \qquad 0\le y(r)\le1.
 \label{eq:certificate-phi-range}
\end{equation}
For the lower inequality, subtract $E(r)$ from $\varphi(r)$ and use
$I_r\ge\Delta E(r)$; for the upper inequality use
$I_r\le\Delta$ and $t+\Delta\le1$.

\paragraph{Elementary operations used by the guards.}
Here and below an intermediate is evaluated in the order prescribed by
the source. Repeated application of
\eqref{eq:certificate-fp-model} gives the following bounds. The table
also records sufficient estimates for each operation, so that no
unquantified error-propagation assertion is needed. Error bounds in its
last column are multiples of $\varepsilon_{\mathrm{in}}$.
\begin{center}
\begin{tabular}{lll}
\hline
Computed intermediate & Exact target & Error bound\\
\hline
$\operatorname{fl}(\widehat b\widehat\Delta)$
 &$b\Delta$ &$4\Delta$\\
$\operatorname{fl}(\operatorname{fl}(\widehat b\widehat\Delta)\texttt{vv[h]})$
 &$b\Delta (r+e^{-r}-1)$ &$54$\\
$\operatorname{fl}(y\widehat b)$ &$yb$ &$9y/8$\\
$\operatorname{fl}(\operatorname{fl}(y\widehat b)\widehat\Delta)$
 &$yb\Delta$ &$4\Delta y$\\
$\operatorname{fl}(\operatorname{fl}(\operatorname{fl}(y\widehat b)
       \widehat\Delta)\texttt{vv[h]})$
 &$yb\Delta (r+e^{-r}-1)$ &$54y$\\
$\operatorname{fl}(\widehat c x)$ &$cx$ &$2y$\\
$\operatorname{fl}(\widehat D-\widehat b)$ &$D-b$ &$3$\\
$\operatorname{fl}(y\operatorname{fl}(\widehat D-\widehat b))$
 &$y(D-b)$ &$4y$\\
\hline
\end{tabular}
\end{center}
For example, the first line follows from
\[
 \Delta\varepsilon_{\mathrm{in}}
 +(2+\varepsilon_{\mathrm{in}})\Delta\varepsilon_{\mathrm{in}}
 +\mathfrak u(2+\varepsilon_{\mathrm{in}})
       \Delta(1+\varepsilon_{\mathrm{in}})+2^{-1075}
 <4\Delta\varepsilon_{\mathrm{in}}.
\]
For the second line the propagated input error and the new rounding
error are at most
\[
 4\Delta\,8192\varepsilon_{\mathrm{in}}
 +0.004\varepsilon_{\mathrm{in}}
 +25.61\mathfrak u+2^{-1075}
 <54\varepsilon_{\mathrm{in}}.
\]
The same calculation multiplied by $y$ proves the fifth line.
The other lines use $|D-b|\le2$, $c<0.9$, and
$\mathfrak u=\varepsilon_{\mathrm{in}}/32$.
The final additions in the code therefore give
\begin{align}
 |\widehat p-p|&<8\varepsilon_{\mathrm{in}}y,
 &p&=cx+y(D-b), \label{eq:certificate-p-error}\\
 |\widehat n(r)-n(r)|&<64\varepsilon_{\mathrm{in}}y,
 &n(r)&=cx+yb\Delta (r+e^{-r}-1), \label{eq:certificate-n-error}
\end{align}
because $|p|\le3y$ and $0\le n(r)\le26.5y$.
The right-hand side of the leaf clipping test is
$\widehat n_{\mathrm{clip}}=\operatorname{fl}(\operatorname{fl}(\widehat c y)
                    +\varepsilon_{\mathrm{ar}}y)$.
Multiplication by $\varepsilon_{\mathrm{ar}}$ is exact, and
\begin{equation}
 |\widehat n_{\mathrm{clip}}-(cy+\varepsilon_{\mathrm{ar}}y)|
 <2\varepsilon_{\mathrm{in}}y.
 \label{eq:certificate-clip-threshold-error}
\end{equation}
The factor $y$ in both bounds is essential for small state coordinates.

For the node guard the code computes
$\widehat\rho=\operatorname{fl}(\operatorname{fl}(x/y)
 +\operatorname{fl}(\widehat\omega/\widehat c))$,
where $\widehat\omega$ is the second intermediate in the table.
Since $b\Delta (r+e^{-r}-1)\le25.6$, $(b\Delta (r+e^{-r}-1))/c<31$, and
$\widehat c>0.84$, division has error at most
\[
 \frac{(54+31)\varepsilon_{\mathrm{in}}}{0.84}
 +32\mathfrak u+2^{-1075}.
\]
Including the division $x/y$ and the final addition yields
\begin{equation}
 |\widehat\rho-\rho(r)|<128\varepsilon_{\mathrm{in}}.
 \label{eq:certificate-ratio-error}
\end{equation}
Likewise, for the product in the upper-endpoint guard,
\begin{equation}
 |\operatorname{fl}(\widehat r\widehat\Delta)-r\Delta|
 <\bigl(0.002+12.8+0.401\bigr)\varepsilon_{\mathrm{in}}
 <16\varepsilon_{\mathrm{in}}.
 \label{eq:certificate-rate-product-error}
\end{equation}

\paragraph{Validity of the zero-continuation and range tests.}
If the leaf test skips interpolation, then $\widehat n(r)\ge\widehat n_{\mathrm{clip}}$.
Equations \eqref{eq:certificate-n-error} and
\eqref{eq:certificate-clip-threshold-error} imply
\[
 n(r)-cy\ge
  (\varepsilon_{\mathrm{ar}}-66\varepsilon_{\mathrm{in}})y>0.
\]
Thus $\rho(r)\ge1$, the exact successor is on the diagonal, and its
continuation is zero. If the leaf test instead evaluates the successor,
the reverse estimates give
\[
 \rho(r)<1+
 \frac{\varepsilon_{\mathrm{ar}}+66\varepsilon_{\mathrm{in}}}{c}
 <1+2\varepsilon_{\mathrm{ar}}.
\]
For the node test, $\widehat\rho\ge1+\varepsilon_{\mathrm{ar}}$ implies
$\rho(r)>1$ by \eqref{eq:certificate-ratio-error}, so all its descendant
continuations are zero by the monotonicity of $\rho$.
Otherwise $\rho(r)<1+2\varepsilon_{\mathrm{ar}}$ again.
In every case where a lower-endpoint successor is evaluated,
$(r+e^{-r}-1)\ge r-1$ now implies
\begin{equation}
 r\Delta\le\Delta+\frac cb\rho(r)
 <\frac1{640}+3.6(1+2\varepsilon_{\mathrm{ar}})<4.
 \label{eq:certificate-successor-rate-range}
\end{equation}
For an upper endpoint, the code evaluates its successor only if
$\operatorname{fl}(\widehat r\widehat\Delta)
 \le4-\varepsilon_{\mathrm{ar}}$.
Equation \eqref{eq:certificate-rate-product-error} proves $r\Delta<4$
in this case as well. Otherwise the node calculation uses the valid
lower bound zero for that successor's $y$-coordinate.

\paragraph{Errors in $\varphi$ and the successor coordinates.}
The following table follows the parenthesization of
\texttt{ph}. Its error bounds are again in units of
$\varepsilon_{\mathrm{in}}$.
\begin{center}
\begin{tabular}{lll}
\hline
Computed intermediate & Exact target & Error bound\\
\hline
$\operatorname{fl}(\operatorname{fl}(1-\widehat t)-\widehat\Delta)$
 &$1-t-\Delta$ &$3$\\
$\operatorname{fl}(\widehat E\widehat K)$
 &$E(1-t-\Delta)$ &$5$\\
$\operatorname{fl}(\widehat t+\widehat I)$ &$t+I_r$ &$4$\\
$\operatorname{fl}(\widehat T+\widehat E\widehat K)$
 &$t+I_r+E(1-t-\Delta)$ &$10$\\
$\operatorname{fl}(y\widehat Z)$ &$y[t+I_r+E(1-t-\Delta)]$ &$11y$\\
\hline
\end{tabular}
\end{center}
In the fourth line $\widehat E\widehat K$ denotes the already rounded
second-line product. Here $\widehat K,\widehat T,\widehat Z$ denote
the corresponding previously computed quantities.
For completeness, the first-line estimate is bounded by
$2+\Delta+(1+\Delta)/32$ in the stated units, with an arbitrarily
small allowance for $\varepsilon_{\mathrm{in}}$ and the subnormal term;
the next lines are bounded respectively by
$3+1+1/32$, $2+1+1/32$, $5+4+1/32$, and
$10y+y/32$, with the same strict slack.
The quantity $1-y$ is an exactly representable multiple of $2^{-30}$.
The last addition thus gives
\begin{equation}
 |\widehat\varphi(r)-\varphi(r)|
 <16\varepsilon_{\mathrm{in}}.
 \label{eq:certificate-phi-error}
\end{equation}

For every evaluated successor, \eqref{eq:certificate-successor-rate-range}
and \eqref{eq:certificate-phi-range} give
$\varphi(r)>e^{-4}>1/64$, hence
$\widehat\varphi(r)>1/65$.
The rounded numerator $\widehat a=\operatorname{fl}(y\widehat E)$
satisfies $|\widehat a-yE|<2\varepsilon_{\mathrm{in}}y$.
Using $y(r)\le1$, the error before rounding the division is at most
\[
 \left|\frac{\widehat a}{\widehat\varphi}-y(r)\right|
 \le\frac{|\widehat a-yE|+y(r)|\widehat\varphi-\varphi|}
             {\widehat\varphi}
 <65(2y+16)\varepsilon_{\mathrm{in}}.
\]
The division adds less than $2\mathfrak u+2^{-1075}$, proving
\begin{equation}
 |\widehat y(r)-y(r)|<2048\varepsilon_{\mathrm{in}}.
 \label{eq:certificate-successor-y-error}
\end{equation}

The leaf computes an \emph{unclipped} $x$-coordinate by dividing
$\operatorname{fl}(\widehat E\widehat n)$ by
$\operatorname{fl}(\widehat c\widehat\varphi)$.
Under its evaluation guard, $n=cy\rho<y$ and
$\widetilde x(r)\le1+2\varepsilon_{\mathrm{ar}}<1.001$.
Equation \eqref{eq:certificate-n-error} and one multiplication give
numerator error less than $66\varepsilon_{\mathrm{in}}y$.
Equation \eqref{eq:certificate-phi-error} gives denominator error less
than $16\varepsilon_{\mathrm{in}}$, and the computed denominator is
greater than $1/80$. The same quotient identity therefore bounds the
error before the final rounding by
\[
 80\bigl(66y+16\cdot1.001\bigr)\varepsilon_{\mathrm{in}}.
\]
Including the division rounding yields
\begin{equation}
 |\widehat x(r)-\widetilde x(r)|
 <8192\varepsilon_{\mathrm{in}}.
 \label{eq:certificate-successor-x-error}
\end{equation}

For a finite upper endpoint that passes its range guard, the node
instead computes
$\widehat x_{\mathrm{low}}
 =\operatorname{fl}(\widehat y(r_+)\min\{1,\widehat\rho(r_-)\})$.
The map $z\mapsto\min\{1,z\}$ is $1$-Lipschitz.
Combining \eqref{eq:certificate-ratio-error} and
\eqref{eq:certificate-successor-y-error}, then allowing one product
rounding, gives
\begin{equation}
 \bigl|\widehat x_{\mathrm{low}}
    -y(r_+)\min\{1,\rho(r_-)\}\bigr|
 <4096\varepsilon_{\mathrm{in}}.
 \label{eq:certificate-node-x-error}
\end{equation}
For example, the coefficient of $\varepsilon_{\mathrm{in}}$ is bounded
by $2048(1+128\varepsilon_{\mathrm{in}})+128+2/32$, which is below
$4096$. When the upper-endpoint guard is not used, the assigned value
is exactly zero and needs no floating-point estimate.

\paragraph{Shifted interpolation and multiplication.}
Let $\widetilde v_{k+1}$ be the ordinary triangular interpolant already
defined above; its vertex values are stored integers divided by $2^{40}$.
The monotonicity passes ensure $\widetilde v_{k+1}\le0$, that
$\widetilde v_{k+1}$ is nonincreasing in its first coordinate and
nondecreasing in its second coordinate, and that it is zero on the diagonal.
The interpolation routine rounds $\widehat y+\varepsilon_{\mathrm{ar}}$
up to a multiple of $2^{-30}$ and
$\widehat x-\varepsilon_{\mathrm{ar}}$ down to such a multiple,
then clips the former to $[0,1]$ and the latter below at zero.
The additions incur less than
$\varepsilon_{\mathrm{in}}/8$ error; multiplication by $2^{30}$ is
exact. Equations \eqref{eq:certificate-successor-y-error}--
\eqref{eq:certificate-node-x-error} and
$\varepsilon_{\mathrm{ar}}=65536\varepsilon_{\mathrm{in}}$
therefore put the integer query above the exact required
$y$-coordinate and below the exact required $x$-coordinate.
Clipping at $0$ or $1$ preserves these inequalities.

For a leaf with $\rho(r)<1$, the required $x$-coordinate is
$\widetilde x(r)=x(r)\le y(r)$, so the queried value is at most
$\widetilde v_{k+1}(y(r),x(r))$. If the routine returns zero because its integer
$x$-coordinate is no smaller than its integer $y$-coordinate, these
coordinate inequalities force the exact successor to be on the
diagonal; the zero return is then valid. For a leaf with $\rho(r)\ge1$,
the true continuation is already zero. In this case either a zero
return or any nonpositive interpolated value is conservative, even
though the code formed the raw coordinate $\widetilde x(r)$ rather
than its clipped version.
The identical argument applies to a node when $\rho(r_-)<1$ and the
required coordinate is $y(r_+)\rho(r_-)$; when $\rho(r_-)\ge1$, every
descendant continuation is zero. These are exactly the cases used by
the exact-arithmetic node bound.

Inside a containing triangle, the code forms nonnegative integer
barycentric numerators with their exact sum equal to the positive
denominator, and takes a downward integer quotient. Thus the returned
integer $v$ satisfies $v/2^{40}\le\widetilde v_{k+1}$ at the integer query. Integer
range bounds are established separately below.

Finally, the code uses
\[
 M=\left\lceil
     \operatorname{fl}\bigl((\widehat\varphi+
                       \varepsilon_{\mathrm{ar}})2^{40}\bigr)
    \right\rceil+1,
\]
where the parenthesized addition is rounded once and multiplication
by $2^{40}$ is exact. Its addition error is less than
$\varepsilon_{\mathrm{in}}/8$. Consequently
\[
 \varphi\le\frac M{2^{40}}
 \le\varphi+2\varepsilon_{\mathrm{ar}}+2\cdot2^{-40}.
\]
Since $v\le0$, the exact downward quotient used for
$Mv/2^{40}$ certifies a lower bound on the exact multiplier times
the required interpolant. This proves the rounding correctness of
both the leaf continuation and the subtree continuation bound used
for pruning.

\subsection{Certification of the running-cost arithmetic}
\label{sec:cost-roundoff}

This subsection concerns the frozen functions \texttt{B}, \texttt{Bnode}, and
\texttt{log\_lower}; it does not introduce a different running-cost bound.
We use the constants $\varepsilon_{\mathrm{in}}=2^{-48}$,
$\varepsilon_{\mathrm{ar}}=2^{-32}$, and $\mathfrak u=2^{-53}$ fixed above.
Write $\mathfrak g_j=j\mathfrak u/(1-j\mathfrak u)$.
Hats denote computed binary64 quantities.  In particular,
$|\widehat p-p|<8\varepsilon_{\mathrm{in}}$ and
$|\widehat t-t|<2\varepsilon_{\mathrm{in}}$, as established in the
continuation analysis.  The validated data satisfy
\[
 \begin{gathered}
 \Delta=1/640,\quad 0\le t\le1-\Delta,\quad 0\le x<y\le1,\\
 \tfrac14<b<2,\quad |p|<2,\quad 0\le I_r\le\Delta,\quad 0\le E(r),e^{-r}\le1.
 \end{gathered}
\]
Here $p=cx+y(D-b)$, where $.85<c<.9$ and $0\le D\le1$: thus
$p>-2$ and $p\le y(.9+1-1/4)\le1.65<2$.
Each stored coefficient has absolute error less than
$\varepsilon_{\mathrm{in}}$, and the stored $\Delta$ has relative error
less than $\varepsilon_{\mathrm{in}}$.  Write
\[
 p_-=-(-p)^+,\qquad p_+=p^+,\qquad
 C(r)=p_-\Delta+p_+I_r,
\]
\[
 \alpha(r)=\left(ybE(r)-\frac{p_-I_r}\Delta\right)\Delta^2,
 \quad \beta(r)=b(1-y+ty)\Delta+ybI_r\Delta,
 \quad z(r)=\beta(r)e^{-r}.
\]

\paragraph{A reproducible error bound for the coefficients.}
For every elementary operation, including one with a subnormal result, we
use the following weaker form of the preceding floating-point model:
\begin{equation}
 |\operatorname{fl}(v)-v|\le\mathfrak u|v|+2^{-1074}.
 \label{eq:cost-basic-roundoff}
\end{equation}
For example, if $|a|\le A$, $|b|\le B$ and the operand errors are
$e_a,e_b$, a product has error at most
\[
 Ae_b+Be_a+e_ae_b+
 \mathfrak u(A+e_a)(B+e_b)+2^{-1074}.
\]
An addition has error at most
$e_a+e_b+\mathfrak u(A+B+e_a+e_b)+2^{-1074}$.
For division by a positive exact $b\ge b_0>e_b$, first bound the input
perturbation by $(e_a+|a/b|e_b)/(b_0-e_b)$, and then apply
\eqref{eq:cost-basic-roundoff}.  The maps $v\mapsto\min(v,0)$ and
$v\mapsto\max(v,0)$ do not increase absolute errors.

The following table applies these inequalities in the source's evaluation
order.  The middle column bounds the magnitude of the corresponding exact
expression, and the last column is a strict error bound in units of
$\varepsilon_{\mathrm{in}}$.  Each multiplication by $\Delta$ is a
separate operation.  The dyadic quantity $1-y$ is exact.
\begin{center}
\small
\begin{tabular}{lll}
\hline
Expression & Exact magnitude bound & Error coefficient\\
\hline
$yb$; $ybE$ & $2$; $2$ & $2$; $5$\\
$p_-I_r$; $p_-I_r/\Delta$ & $2\Delta$; $2$ & $3$; $2000$\\
$ybE-p_-I_r/\Delta$ & $4$ & $2010$\\
$(ybE-p_-I_r/\Delta)\Delta$ & $4\Delta$ & $4$\\
$\alpha=(ybE-p_-I_r/\Delta)\Delta^2$ & $4\Delta^2$ & $1$\\
$ty$; $1-y+ty$ & $1$; $1$ & $3$; $4$\\
$b(1-y+ty)$; $b(1-y+ty)\Delta$ & $2$; $2\Delta$ & $10$; $1$\\
$ybI_r$; $ybI_r\Delta$ & $2\Delta$; $2\Delta^2$ & $3$; $1$\\
$\beta$; $z=\beta e^{-r}$ & $2\Delta(1+\Delta)<.004$ & $3$; $4$\\
$p_-\Delta$; $p_+I_r$; $C=p_-\Delta+p_+I_r$ & $2\Delta$ & $1$; $3$; $5$\\
\hline
\end{tabular}
\end{center}
For the largest intermediate error, for example, division of the preceding
$3\varepsilon_{\mathrm{in}}$ error by
$\Delta(1-\varepsilon_{\mathrm{in}})$ contributes less than
$1921\varepsilon_{\mathrm{in}}$; denominator perturbation and rounding
contribute less than $3\varepsilon_{\mathrm{in}}$.  The two following
multiplications reduce the bound to the displayed $4$ and $1$.
All other entries follow by the same displayed product and sum rules;
the strict margins also exceed the sum of the fewer than $100$ possible
$2^{-1074}$ terms.  In particular, all four final errors are less than
$8\varepsilon_{\mathrm{in}}$.

\paragraph{The logarithm computed by the frozen source.}
This paragraph certifies the actual twelve-term loop, not the unused
degree-five polynomial mentioned in an older explanation.  Whenever the
logarithmic branch is reached, put
$v=\widehat\beta/\widehat\alpha$ in exact real arithmetic and
$\widehat v=\operatorname{fl}(v)$.  The guards give
$\widehat\beta>\widehat\alpha>\varepsilon_{\mathrm{ar}}$; hence
$1\le\widehat v<2^{25}$ and $|\log\widehat v-\log v|<2\mathfrak u$.
The exact binary operation \texttt{frexp} gives $\widehat v=2^e m$ with
$0\le e\le24$ and $1\le m<2$.  For $w_0=(m-1)/(m+1)$,
\[
 S_{12}(m)=2\sum_{j=0}^{11}\frac{w_0^{2j+1}}{2j+1}
 \le\log m,\qquad 0\le w_0<\tfrac13.
\]
The subtraction $m-1$ is exact.  Rounding of $m+1$ and the division give
a relative error bounded by $\mathfrak g_2$
in a nonzero $w_0$.  In the largest monomial, its degree is $23$;
the repeated squaring-factor and product operations contribute at most
$22$ further factors, division contributes one, and summation contributes
at most eleven.  Thus the computed $2\,\texttt{sum}$ is at most
\[
 (1+\mathfrak g_{128})S_{12}(m)+256\cdot2^{-1074}.
\]
This bound also covers underflow of the smallest positive monomials.
There is no assumption about a platform logarithm implementation.

The hexadecimal constant in the source is the exact rational
\[
 c_2=\frac{3121657384082679}{2^{52}}
       =\texttt{0x1.62e42fefa39eep-1}.
\]
The finite rational inequality
\[
 2\sum_{j=0}^{39}\frac1{(2j+1)3^{2j+1}}
       > c_2+2^{-53}
\]
is checked by multiplication by the positive common denominator.
The left side is a lower partial sum of the positive series for $\log2$,
so $c_2<\log2$.  Since $S_{12}(m)<1$ and $e\le24$, the above
polynomial error, the product $ec_2$, the final addition and subtraction,
and the rounded quotient have combined upper error at most
\[
 \mathfrak g_{128}S_{12}(m)+24\mathfrak u+25\mathfrak u
 +26\mathfrak u+2\mathfrak u+512\cdot2^{-1074}
 <256\mathfrak u+512\cdot2^{-1074}
 <32\varepsilon_{\mathrm{in}}.
\]  Consequently the returned value satisfies
\begin{equation}
 \texttt{log\_lower}(\widehat v)
 \le\log v-\varepsilon_{\mathrm{ar}}
            +32\varepsilon_{\mathrm{in}}
 <\log v.
 \label{eq:cost-log-certified}
\end{equation}

\paragraph{The free bound and its branch tests.}
Recall $\mathcal H$ from~\eqref{eq:cert-iid-free}.
On the guarded domain and on the segment joining the exact and computed
coefficient pairs, $\alpha>\varepsilon_{\mathrm{ar}}/2$ and
$\beta<.004$.  The function is continuously differentiable across
$\beta=\alpha$, and direct differentiation gives
\[
 0\le\partial_\beta \mathcal H\le1,\qquad
 0\le\partial_\alpha \mathcal H<10.
\]
Indeed $\beta/\alpha<2^{26}$ in the logarithmic region.  Thus coefficient
perturbations change $\mathcal H$ by less than
$88\varepsilon_{\mathrm{in}}$.  This argument permits the computed and
exact pairs to fall on different sides of the branch test.
In the rational branch the division, subtraction, and multiplication
contribute less than $\varepsilon_{\mathrm{in}}$ in absolute error.
In the other branch, \eqref{eq:cost-log-certified} and the remaining two
rounded operations give the same one-sided bound.
Including the error of \texttt{base}, addition, and subtraction, the
computed candidate \texttt{base+freebound-eps} is therefore at most
\[
 C(r)+\mathcal H(\alpha(r),\beta(r))
       -\varepsilon_{\mathrm{ar}}+128\varepsilon_{\mathrm{in}},
\]
and is a valid lower bound.  If a positivity guard fails, the source simply
does not use this candidate.  The baseline \texttt{min(p,0)*dt-eps}
is separately a lower bound for $\min(p,0)\Delta$.

\paragraph{The scalar shifted bound.}
The frozen rate grid contains $64$ exactly and every other rate has
distance at least $16$ from $64$.  Thus \texttt{rr[hi]<=64} cannot admit
an exact rate exceeding $64$.  The test
$\widehat z>\operatorname{fl}(\widehat\alpha+
\varepsilon_{\mathrm{ar}})$ implies $z>\alpha$: the total coefficient
and addition errors are less than $17\varepsilon_{\mathrm{in}}$.
The same strict inequality holds throughout the segment between the
exact and computed $(\alpha,z)$ pairs.  For
\[
 \mathcal J(\alpha,z,r)=z+\alpha r/2-
       \frac{\alpha^2}{24(z-\alpha/2)},\qquad z>\alpha\ge0,
\]
differentiation gives
\[
 1\le\partial_z\mathcal J\le\tfrac76,
 \qquad |\partial_\alpha\mathcal J|\le33,
 \qquad |\partial_r\mathcal J|=\alpha/2<1
 \quad(0\le r\le64+\varepsilon_{\mathrm{in}}).
\]
Hence coefficient perturbations, including that of $r$, contribute less
than $275\varepsilon_{\mathrm{in}}$.  The subtraction in
$\ell=z-\alpha/2$ has no harmful relative cancellation, since
$\ell>z/2$ and $\ell>\alpha/2$.
The exact quotient $\alpha^2/(24\ell)$ is at most $\alpha/12$.
The rounding of $\ell$, its multiplication by $24$, the numerator
product, and division give absolute error at most
$\mathfrak g_5\alpha/12+2^{-1000}$ in this quotient.
For the additive term, the guard gives $\ell>\varepsilon_{\mathrm{ar}}/2$;
propagating the fewer than ten possible $2^{-1074}$ errors through these
operations amplifies them by less than $2^{40}$, which is covered by
$2^{-1000}$.  All remaining products and sums have magnitude less
than $1$; applying \eqref{eq:cost-basic-roundoff} to them and to
\texttt{base} bounds the complete upward error by
$512\varepsilon_{\mathrm{in}}$.
Thus \texttt{sharp-eps} is below the exact shifted expression whenever
this branch is used.

\paragraph{The shifted bound for a tree node.}
Let the node's endpoint rates be $r_-\le r_+\le64$.
The quantities actually formed by the source are
\[
 \overline\alpha=\operatorname{fl}(\widehat\alpha(r_-)
                 +16\varepsilon_{\mathrm{in}}),\quad
 \underline\alpha=\max\{0,\operatorname{fl}(\widehat\alpha(r_+)
                 -16\varepsilon_{\mathrm{in}})\},
\]
\[
 \underline z=\max\{0,\operatorname{fl}(\widehat z(r_+)
                 -16\varepsilon_{\mathrm{in}})\},\quad
 \underline r=\max\{0,\operatorname{fl}(\widehat r_-
                 -16\varepsilon_{\mathrm{in}})\}.
\]
The table above proves
$\overline\alpha\ge\alpha(r_-)$,
$0\le\underline\alpha\le\alpha(r_+)$, and
$0\le\underline z\le z(r_+)$: the extra additions and subtractions
have error less than $\varepsilon_{\mathrm{in}}$.
For the last expression their error is at most
$3\varepsilon_{\mathrm{in}}$, so $0\le\underline r\le r_-$ as well.
The rounded guard implies $\underline z>\overline\alpha$.
Therefore the exact expression in these already directed floating inputs,
\[
 C(r_+)+\underline z+
 \underline\alpha\,\underline r/2-
 \frac{\overline\alpha^2}
      {24(\underline z-\overline\alpha/2)},
\]
is a lower bound for the shifted cost of every descendant original row,
by the endpoint argument proved above.
Its quotient is again at most $\overline\alpha/12$, with absolute
operation error at most $\mathfrak g_5\overline\alpha/12+2^{-1000}$.  The remaining operation errors and
the error of \texttt{base} together are less than
$16\varepsilon_{\mathrm{in}}$.  Subtracting
$2\varepsilon_{\mathrm{ar}}$ therefore preserves the lower-bound
direction.  No monotonicity of the full shifted expression is assumed.

In \texttt{B}, the maximum of the available certified candidates remains
certified.  The final subtraction of $\varepsilon_{\mathrm{ar}}$, exact
power-of-two scaling, floor, and subtraction of one integer unit give a
lower bound in units of $2^{-40}$.  The analogous floor and subtraction
in \texttt{Bnode} have the same direction.  The magnitudes just established
are below $1$, so the final floating subtraction has error less than
$\varepsilon_{\mathrm{in}}$ and cannot exhaust any of these margins.

\subsection{Table magnitudes, integer safety, and the final objective}
\label{sec:integer-safety}

The allowance computation requires no approximation theorem.  The
coordinates are exact multiples of $2^{-30}$, so its $w=y-x$ is exact
and lies in $[0,1]$.  Whatever error occurs in computing the intermediate
\texttt{future}, the two comparisons in
\texttt{max(0,min(w,w-future))} put the returned floating value in $[0,w]$.
Multiplication by $2^{24}$ is exact, and consequently
\[
 0\le\texttt{allowance}\le2^{24}.
\]
Subtracting this integer therefore only lowers a proposed table value,
by at most $2^{-16}$.  The fixed additional slack is $2^{-28}$.

Let $M_k$ be the largest absolute table entry in layer $k$, after division
by $2^{40}$.  All entries are nonpositive.  Integer barycentric
interpolation of the next layer lies between its minimum and zero:
flooring cannot cross the integer minimum.  The previously certified
multiplier, including its ceiling and extra integer unit, is at most
$2^{40}(1+2\varepsilon_{\mathrm{ar}}+2\cdot2^{-40})$.
The subsequent integer floor loses less than $2^{-40}$.
For the running cost, the actual baseline arithmetic gives
\[
 2^{-40}\texttt{B(hi)}
 \ge -3/640-3\varepsilon_{\mathrm{ar}}-2\cdot2^{-40}.
\]
Indeed $|\widehat p|<2+8\varepsilon_{\mathrm{in}}$;
the computed product with $\widehat\Delta$ is greater than $-3/640$;
the baseline and final floating subtractions use two
$\varepsilon_{\mathrm{ar}}$ deductions, their rounding is covered by
the third, and flooring followed by subtraction of one costs less than
two integer units.  Taking a maximum can only improve this lower bound.
It follows, with deliberately enlarged constants, that
\begin{align*}
 M_k&\le aM_{k+1}+d,\qquad M_{640}=0,\\
 a&=1+2\varepsilon_{\mathrm{ar}}+2\cdot2^{-40},\\
 d&=3/640+4\varepsilon_{\mathrm{ar}}+8\cdot2^{-40}
                  +2^{-28}+2^{-16}.
\end{align*}
Taking minima with zero or with additional proposals, and either
monotonicity-lowering pass, preserves this global lower bound.
Since $640(a-1)<1$, the binomial theorem gives
\[
 M_k\le d\sum_{j=0}^{639}a^j
 \le\frac{640d}{1-640(a-1)}
 =\frac{1846693964}{613566573}<\frac{301}{100}<4.
\]
Thus every stored integer has magnitude below $4\cdot2^{40}$.

Whenever interpolation is performed, its tests give
$0\le X<Y\le2^{30}$.  Lookup indices lie in $[0,65536]$, and every
grid width and offset lies in $[0,2^{30}]$.  The barycentric denominator
and each unscaled weight are at most $2^{60}$, so their products and
differences fit signed 64-bit integers.  The selected triangle has
nonnegative weights summing to the denominator.  Its weighted numerator
and all partial sums therefore have magnitude at most
$4\cdot2^{40}\cdot2^{60}=2^{102}$, safely within signed 128-bit range.
Continuation multiplier products are smaller than $2^{84}$, and the
negative-floor formula's additional denominator is also harmless at
these bounds.  Each conversion of a running-cost value, multiplier, or
allowance to a signed 64-bit integer is within its range.
The state-array size is $1201\cdot1202/2=721801$; all index products fit
signed 32-bit integers.  A binary tree over $131093$ leaves has fewer
than twice that many nodes; the additional proposal evaluations are
fewer than $50$ per state.  Per-state counters fit unsigned 32-bit
integers and their sums over all states fit unsigned 64-bit integers.

Finally, input validation establishes
$\texttt{HU}=\lfloor2^{40}\theta(0)\rfloor$.
The reported exact objective is therefore
\[
 \frac{\texttt{HU}+\texttt{table}[m,0]}{2^{40}}
 \le\theta(0)+\frac{\texttt{table}[m,0]}{2^{40}}.
\]
Together with the original-row feasibility induction, this proves the
claimed direction of the rational objective; the printed decimal plays
no role in acceptance.


\section{Source mapping and reproduction}
\label{sec:reproduction}
The source lines cited in this document refer to the unchanged
\texttt{certificate/verify.cpp}. Its SHA-256 is
\begin{center}
\small\texttt{6cafb1fb9e1be830c8b2c8cd60304bb1e2c3496e3da5a3e7e2d47cb4b9a84219}.
\end{center}
The release preflight checks this identity and all 41 entries in the
current certificate manifest. In particular, the mathematical argument is tied to
the actual expression order and guards of this source.

\begin{longtable}{@{}p{0.25\linewidth}p{0.21\linewidth}p{0.46\linewidth}@{}}
\toprule
Source item & Paper notation & Role\\
\midrule
\endhead
\texttt{N}, \texttt{Q}, \texttt{m} & $N_t,N_q,N_y$ & Time, rate, and state grid dimensions\\
\texttt{bs[k]}, \texttt{bl[k]} & $b$ on $[t_k,t_{k+1}]$ & Identical approximations of the same exact rational drop rate\\
\texttt{DD[k]} & $1-\theta(t_k)$ & Frozen threshold complement\\
\texttt{cap} & $c$ & The constant from Appendix A.1\\
\texttt{ee[h]}, \texttt{ii[h]} & $e^{-\Delta r_h},I_{r_h}$ & Rate-dependent input coefficients\\
\texttt{vv[h]}, \texttt{er[h]} & $r_h+e^{-r_h}-1,e^{-r_h}$ & Rate-dependent input coefficients\\
\texttt{A}, lines 63--68 & Continuation in \eqref{eq:cert-iid-row} & Leaf continuation with conservative coordinates and multiplier\\
\texttt{B}, lines 78--96 & Lemma A.3 cost & Conservative evaluation of available cost bounds\\
\texttt{rectA}, lines 98--106 & Lemma \ref{lem:cert-iid-continuation} & Continuation bound for a whole rate subtree\\
\texttt{Bnode}, lines 108--124 & Lemmas \ref{lem:cert-iid-free}--\ref{lem:cert-iid-shifted} & Cost bound for a whole rate subtree\\
\texttt{search}/\texttt{solve}, lines 135--155 & Theorem \ref{thm:cert-iid-backward} & Complete constraint coverage and downward candidate updates\\
\texttt{lower\_monotone}, lines 158--166 & Lemma \ref{lem:cert-iid-lowering} & Current-layer lowering and explicit monotonicity checks\\
\bottomrule
\end{longtable}

The scalar storage constant named \texttt{U} in the source is $2^{40}$,
and \texttt{L} is the state-grid denominator $2^{30}$. These code
identifiers do not redefine the paper's mathematical symbols.

\subsection{Input validation}
The input validator represents every hexadecimal input exactly
as a rational number. For $0\le z\le1$, the alternating Taylor sums of
degrees 25 and 24 enclose $e^{-z}$. For larger arguments below 64,
the validator repeatedly halves the argument until it lies in $[0,1]$
and squares the positive rational endpoints to recover the original
exponential. For $z\ge64$, it uses
\[
 0\le e^{-z}\le 2^{-\lfloor z\rfloor},
\]
which follows from $e>2$. These enclosures validate the entire frozen
coefficient file, including the differences used to form $I_r$ and
$r+e^{-r}-1$, with error strictly below
$\varepsilon_{\mathrm{in}}$. The exact $r=0$ value is $I_0=\Delta$.
The extra preflight checks establish the fixed-instance facts used by
the arithmetic proof and the consistency of all input representations.

\subsection{Commands and acceptance conditions}
From the release root, run
\begin{verbatim}
python3 reproduce.py --check-only
python3 reproduce.py --full --threads 8
\end{verbatim}
The first command checks the current certificate manifest, the exact preflight
assumptions, and all coefficient input enclosures. The second also invokes
the launcher, reconstructs all 640 backward layers, and checks
the exact objective against \texttt{certificate/RESULTS.json}. It compares
the complete final current table, next-layer table, and hint array with
the original checkpoint hashes. Acceptance never uses a decimal display.

The launcher requires Python 3.9 or later and a GNU-compatible C++17
compiler with OpenMP and signed \texttt{\_\_int128}. The compilation
options include \texttt{-ffp-contract=off}; fast-math and disabled
assertions are rejected. The evaluator sets and checks round-to-nearest
on each worker before its state computations. Python assertions must
also remain enabled; the reproduction script, launcher, coefficient
validator, and row checkers reject disabled Python assertions. Standard IEEE gradual underflow is required, with flush-to-zero
and denormals-are-zero modes disabled.

The files in \texttt{validation/} identify the checks performed. The dated
complete replay record and layer log concern all 640 layers of the unchanged
evaluator and inputs; current packaging checks are recorded separately.
Selected-state diagnostic scans and exact interval row checks are
additional, explicitly scoped evidence. Their role is to cross-check
the implementation; the full feasibility argument is the induction and
arithmetic proof above.

The source, rational inputs, and proof suffice to verify the fixed
threshold. No property of the historical threshold-search procedure is
used. To rebuild this PDF, run \texttt{latexmk -pdf verification.tex}
in \texttt{docs/}; it repeats LaTeX until references and the table of contents
stabilize. The distributed PDF is already compiled.

\end{document}
